\section{The exact contraction interface used as input}\label{sec:interface}

The graph is bipartite and Boolean: every edge carries one bit. A vertex
entry is specified by a uniform procedure. With $E$ the edge set, the
contraction is
\[
 Z=\sum_{x\in\{0,1\}^{E}}\prod_v T_v(x|_{\delta(v)}).
\]
There are at most $s$ edges and polynomially many vertices. Entries belong
to $(\Z/2^{P_0}\Z)[V]/(V^\ell)$, with $1\le P_0,\ell\le s$.
The description length is $\LL$, and an original entry costs
$(s+\LL)^{O(1)}2^{O(s^{9/10})}$ time and polynomial space. The
constants in these premises are fixed. Clause and equality tensors are
much cheaper than this allowance.

The unchanged algebraic ingredients are the Cauchy deformation and
Hermite recovery of \cite{source}, the sharp precision and near-unit
packing refinements \cite{precision,asymmetric}, and the full-block
integral Cauchy modes of \cite{joint}. They operate over
\[
 R_p=(\Z/2^p\Z)[U,V]/(f(U),V^\ell),
\]
where $f$ is monic and irreducible modulo two, with degree polynomial
in $s$. Its degree, $\ell$, the precisions, and coefficient descriptions
remain polynomially bounded. A tensor retains its creation precision:
later calls use the same canonical representatives of its creation-time
entries. Canonical lifting is not assumed to preserve ring operations.

\subsection{The imported local identities and entropy estimate}

At a call with $M$ wires, fixed group alphabet $r$, taper $u$, packet
exponent $a$, and retained valuation $i<u$, put $D=2^{a-i}$. The cited
lemmas give child size $(2/r+o(1))M$, working precision
\[
 p_t\le(6+o(1))MD/u,
\]
and at most $(2/r^2+o(1))M/D$ mode families of each index on each
directed pair. Their usable weights are
$a+1-i+4j-2\pc(j)$. For an exact weight allocation, let
$\mathcal I_{gh}$ be the finite mode list on directed pair $(g,h)$;
empty lists are skipped. Set
\[
 h_{gh}=\log_2|\mathcal I_{gh}|,
 \qquad \K(M,u)=\kappa\frac M{\sqrt u},
 \qquad\kappa=\frac{\pi\sqrt2}{\ln2}.
\]
The contribution is the exact finite sum
\begin{equation}\label{eq:exact}
 \sum_{\xi_{gh}\in\mathcal I_{gh}}
       \prod_{g=1}^{r} Z_g((\xi_e)_{e\ni g})\pmod {2^{p_t}}.
\end{equation}
Every $Z_g$ entry is an actual recursive child contraction. It is not a
free table lookup until its entry has been computed and stored.

For every vector $z\ge0$ in a fixed finite support family, the imported
common-tilt entropy bound is
\begin{equation}\label{eq:entropy}
 \sum_{g\ne h}z_{gh}h_{gh}
 \le\left(\frac{\norm z}{r}+o(1)\right)\K(M,u).
\end{equation}
For unordered edges set $g_{\{a,b\}}=h_{ab}+h_{ba}$. Thus
\begin{equation}\label{eq:unordered}
 v\cdot g\le
 \left(\frac{\sqrt2\norm v}{r}+o(1)\right)\K(M,u).
\end{equation}
The factor $\sqrt2$ must be retained. The generating-product logarithm
behind this estimate has leading term $\pi^2/(24\theta)$; minimizing
the precision term plus its common-tilt contribution gives
\eqref{eq:entropy}. The additive precision buffer used below makes the
error uniform over the finite support family and all active depths.

There are polynomially many weight allocations, since $r$ is fixed.
The list cardinalities, iteration, and ranking are obtained by the
existing finite knapsack procedure, not by an additional counting oracle.

\subsection{Packing and the fully charged recurrence}

Write $L_p=2^a$, $Q=r^a$, and let $R,C$ be the two side-degree bounds.
The sufficient packing conditions from \cite{asymmetric} are
\begin{equation}\label{eq:gates}
 QR/M=o(1),\qquad Q^2C(\log s)^2/M=o(1),\qquad L_pQ^2/M=o(1).
\end{equation}
The notation $L_p$ distinguishes packet length from the SAT literal count.
At the root both degree bounds are original. Below it the fresh side has
degree at most $2^{a_{d-1}+1}$, and the retained side contains original
tensors and earlier packets; no retained vertex gains ports.

Caching one singleton group table requires at most
$2^{(\alpha_r+o(1))\K}$ paid recursive evaluations, where
\[
 \alpha_r=\frac{\sqrt{2(r-1)}}r.
\]
If local table combination has coefficient $\beta_r$, the recurrence is
\begin{equation}\label{eq:paid}
 T_d\le\poly(s)\left(
 2^{(\alpha_{r_d}+o(1))\K_d}T_{d+1}
 +2^{(\beta_{r_d}+o(1))\K_d}\right).
\end{equation}
The maximum over all routes in this recurrence determines the global
coefficient. In particular, improving only a local matrix bound need not
improve the complete algorithm. We explicitly compute both the local
routes and the terminal route in Section~\ref{sec:schedules}.

